% =====================================================================
% Seção "Conclusão" — RASCUNHO para o autor REESCREVER na própria voz.
% Inserir APÓS a Seção 5 (Resultados/Discussão) e ANTES de "Uso de IA".
% ATENÇÃO: não submeta como está — reescreva com suas palavras (política
% de IA + checagem de plágio). Confira também os números com o seu texto.
% =====================================================================

\section{Conclusão}
\label{sec:conclusao}

Este trabalho investigou, sob a ótica da Engenharia de Software, a migração do
InterpreteLabBR --- um sistema de apoio à decisão clínica para interpretação de
hemogramas fundamentado nos valores de referência da PNS --- do paradigma de
aplicação web progressiva (PWA) para um aplicativo móvel dedicado em React Native.
A avaliação organizou-se em três frentes complementares: a caracterização da
engenharia da migração (Frente~B), as medições de desempenho em dispositivo físico
via ADB (Frente~A) e a inspeção heurística de usabilidade (Frente~C).

Os resultados sustentam a hipótese central de reuso arquitetural. A estratégia
\emph{API-first} permitiu reaproveitar integralmente (100\%) a lógica clinicamente
sensível do \emph{backend}, sem reescrita, confinando o esforço de migração à
camada de apresentação; adicionalmente, a transição eliminou 1.239 linhas de código
de contorno de \emph{WebView} --- cerca de 49\% do código do cliente web ---,
evidenciando, no plano do código-fonte, o \emph{overhead} que a literatura atribui
a esse paradigma. Quanto ao desempenho, as medições indicaram vantagem expressiva
do cliente nativo, com redução de 63,9\% no consumo de memória e rolagem fluida
(0,2\% de quadros com atraso), confirmando parcialmente a hipótese de ganho de
desempenho. A inspeção heurística, conduzida como estudo-piloto, não revelou
barreiras graves de usabilidade e produziu achados qualitativos acionáveis,
sobretudo na tela de apresentação de resultados.

Como contribuição, este trabalho oferece um estudo de caso empírico e reprodutível
de migração de paradigma em um sistema real de saúde pública brasileira, conduzido
com ferramentas gratuitas e metodologia verificável, que pode servir de referência
para transições arquiteturais semelhantes no contexto da saúde digital nacional.
Entre as limitações, destacam-se a medição de desempenho em um único dispositivo de
gama intermediária, o caráter retrospectivo do mapeamento de migração e o painel
reduzido da inspeção de usabilidade, o que confere aos achados dessa frente natureza
exploratória.

Como trabalhos futuros, propõe-se ampliar a inspeção de usabilidade para um painel
maior e incluir as métricas de rolagem ainda pendentes (FPS, CPU e consumo
energético), em dispositivos representativos do perfil de \emph{hardware} dos
usuários do SUS; estender o suporte de extração a leiautes de laudo da rede privada
e à população pediátrica; e evoluir o sistema de uma arquitetura baseada em regras
para uma abordagem híbrida com aprendizado de máquina, mantendo os valores de
referência da PNS como padrão-ouro de rotulagem.
